\section{Additional structure of frozen-ideal protection}\label{anc:section}
The results below retain structural information from the sharp-subsidy
program that is not implied by the numerical bound alone. The compact-core
representation and the finite cyclic example concern sparse supports;
the incompatibility and perturbation results are all-size statements.
They remain parts of the same protected-ideal question.

\subsection{A compact entire-core representation on sparse supports}
Suppose the three pair matrices have support contained in the pairs of
an optimal diagonal matching $\mu$ on $m$ labels and an ordinary
player-disjoint partial matching $P$ of $r$ full teams. A pair is called
supported if it appears in $\mu$ or $P$, even when its value is zero.
Teams of $P$ may repeat diagonal labels or share pairs with $\mu$.
Let $T$ be the diagonal matching's worth and assume the entire core is
nonempty.

\begin{proposition}[Sector-minimum representation]\label{anc:compression}
A vector $x$ is in the entire original core if and only if there exist
three variables $\ell_A,\ell_B,\ell_C$ such that
\begin{enumerate}
\item $x\ge0$, $\ell_s\ge0$, and $\ell_s\le x_i$ for every sector-$s$
player $i$;
\item for each supported pair $\{i,j\}$ of value $a_{ij}$, whose missing
sector is $u$, $x_i+x_j+\ell_u\ge a_{ij}$;
\item the full-team inequalities hold for every team of $\mu\cup P$;
\item for each supported off-diagonal pair joining labels $a\ne b$, its
two full-team inequalities obtained by completing it at label $a$ or
label $b$ hold;
\item $x(N)=T$.
\end{enumerate}
The system uses $3m+3$ variables, at most $7m+10r$ inequalities apart
from variable nonnegativity, and one total equality. No equality forcing
$\ell_s$ to be a coordinate minimum is needed.
\end{proposition}
\begin{proof}
Given a core point, take the actual sector minima for $\ell$.
Completing a supported pair with a player attaining the missing-sector
minimum gives the pair inequality, because all other pair values are
nonnegative. The listed full-team constraints and efficiency also hold.

Conversely, consider an arbitrary full team. If it has two supported
off-diagonal pairs, their common player belongs to one team of the
actual partial matching $P$, so both pairs belong to that same team.
Their union is therefore a team of $P$ already listed. If it has one
supported off-diagonal pair, a third label equal to one endpoint label
gives a listed local completion. Otherwise there is no diagonal pair,
the remaining off-diagonal pairs are unsupported, and its worth equals
that single pair value. The pair inequality and $\ell_u\le x_k$
then give its full-team inequality. With no supported off-diagonal
pair, the team is either diagonal, or has at most one supported diagonal
pair, or has worth zero. The listed diagonal inequality, pair-plus-minimum
inequality, or nonnegativity covers these cases. Summing over disjoint
full teams proves every coalition constraint.

There are $3m$ lower-envelope inequalities, at most $3m+3r$ pair
inequalities, $m+r$ diagonal or $P$ team inequalities, and at most
$6r$ completion inequalities. Their total is $7m+10r$; duplicates can
only reduce it.
\end{proof}
For example, a supported off-diagonal $AB$ pair with labels $a\ne b$
gives, by comparing its two completions with tight diagonals,
\begin{align*}
 x_B(b)-x_B(a)&\ge AB[a,b]+BC[b,a]-AB[a,a]-BC[a,a],\\
 x_A(a)-x_A(b)&\ge AB[a,b]+AC[a,b]-AB[b,b]-AC[b,b].
\end{align*}
The cross terms must be retained for repeated-label $P$ teams. Any valid
core inequality $x_i\ge x_j+c$ also yields $H_i\ge H_j+c$ by using
a core point attaining $H_j$. A cycle of such difference inequalities
has total right-hand side at most zero; if the total is zero, every
inequality is tight at every core point. These are useful consequences,
but discarding the remaining team constraints need not describe the
entire core or its ideals. The compressed system is not claimed for
arbitrary dense supports.

\subsection{Linearly many incompatible ideal faces}
\begin{proposition}\label{anc:attainment}
In the family of Proposition~\ref{prop:family}, let $U$ be the players
outside its $\ell$ special teams. Any original core allocation attains
$H_i=3$ for at most one $i\in U$. Exactly $3(\ell+1)$ core allocations
are necessary and sufficient to attain all those unmatched ideals
collectively.
\end{proposition}
\begin{proof}
There are $\ell+1$ unmatched players in each sector. If player $(s,j)$
receives three at a core point, diagonal tightness makes the other two
coordinates at label $j$ zero. Complete their opposite-sector diagonal
pair of value one with any sector-$s$ player. Its team inequality forces
every sector-$s$ payoff to be at least one.
A second ideal three in another sector would likewise force every
payoff in that sector to be at least one, contradicting its zero payoff
at $j$. If two unmatched ideals in the same sector were attained, that
sector's $\ell+1$ unmatched players would receive at least
$(\ell+1)+4=\ell+5$ in total. But summing the $\ell$ disjoint
special-team constraints gives
\[
                  x(U)\le T-5\ell=\ell+3,
\]
a contradiction. Each unmatched ideal face is nonempty by the witnesses
in Proposition~\ref{prop:family}; choosing one witness per face attains
the lower bound $|U|=3(\ell+1)$ on the required number of points.
\end{proof}
A bounded collection of allocations that exactly attain the ideals
cannot therefore work uniformly over this family. This does not exclude
fractional protection obtained from nonmaximizing allocations. In
contrast, if one core point simultaneously attains every ideal of each
sector, the average of the three points already protects $H/3$ with
zero subsidy. Published side-optimality in the square supermodular class
with local monotonicity~\cite[Theorem~3.2]{Martinez2026} supplies such a
safe class under that theorem's additional hypotheses; neither the
side-optimality concept nor this elementary averaging observation is
claimed as a new contribution.

\subsection{A ring and the effect of weakening one pair}
For $b\ge2$, take $m=4b$ labels and initialize all three pair matrices
to the identity. Let the actual partial matching $P$ consist of
\[
 e_j=(j,(j+1)\bmod b,b+j),\qquad
 f_j=(2b+2j,2b+2j+1,j),\qquad 0\le j<b.
\]
For every displayed team set its off-diagonal pair values, in order
$AB,AC,BC$, to $(2,1,2)$; other off-diagonal values are zero.
Call the resulting market $R_b$. Call $R'_b$ the market obtained by
lowering the single entry $AB[0,1]$ from two to zero. The teams in $P$
are player-disjoint and have three distinct labels.

\begin{theorem}[An all-size perturbation family]\label{anc:ring}
Both $R_b$ and $R'_b$ have nonempty original cores and grand worth
$T=12b$. Their protected subsidies are
\[
                          s(R_b)=0,\qquad s(R'_b)=4b-2.
\]
Every ideal of $R'_b$ equals three. Each market's ideals are calculated
in its own original core and then frozen during protection.
\end{theorem}
\begin{proof}[The unmodified ring]
At each hub label $j<b$ pay $4/3,4/3,1/3$ in sectors $A,B,C$.
At each other label pay $7/3$ to its unique $P$-used player and $1/3$
to its other two players. Denote this vector by $y$. Every diagonal
team receives three, every $P$ team receives five, and $y(N)=12b$.
Two positive off-diagonal pairs in a team force it to be a $P$ team,
because $P$ is a partial matching. Any other nondiagonal team has at
most one such pair and at most one diagonal pair. For a pair of value
two, its endpoints receive at least $8/3$ and its third player at least
$1/3$, covering total worth at most three. For a pair of value one,
the endpoints receive at least $8/3$, covering total worth at most two.
With no off-diagonal pair, worth is at most one, again covered. Hence
$y$ covers every team and every coalition. Together with the diagonal
matching this proves grand worth $12b$ and core nonemptiness.

At any original core point, diagonal tightness and the completion
inequalities for $(j,j+1,j)$ and $(j,j+1,j+1)$, with hub indices taken
modulo $b$, give $x_B(j+1)\ge x_B(j)$ and $x_A(j)\ge x_A(j+1)$.
Around the ring these become equalities. Write the common hub payoffs
as $A,B,C$ with $A+B+C=3$, and put $c=\min_k x_C(k)$.
The ring-team inequality gives $x_C(b+j)\ge2+C$.
Completing the diagonal $AB$ pair at $b+j$ with a sector-$C$ minimum
gives
\[
       1-c\le x_A(b+j)+x_B(b+j)=3-x_C(b+j)\le1-C.
\]
Thus $c=C$, $x_C(b+j)=2+c$, and the other two payoffs there sum to
$1-c$. In particular $0\le c\le1$ and those two coordinates are at
most one.

For attached labels $u=2b+2j$, $v=2b+2j+1$, let
$a=\min_kx_A(k)$ and $d=\min_kx_B(k)$. With $p=x_A(u)$ and
$q=x_B(v)$, the attached-team inequality is $p+q+c\ge5$.
Diagonal tightness gives $q\le3-a-c$ and $p\le3-d-c$.
Hence $p\ge2+a$, $q\ge2+d$, and each of the other coordinates at
$u,v$ is at most one. All these estimates hold on the entire core.
Thus unmatched players and hub-$C$ players have ideal at most one;
all other players have ideal at most three by diagonal tightness.
Vector $y$ pays the former group $1/3$ and every other player at least
one, so it is already an efficient protected allocation.
\end{proof}

\begin{proof}[The broken ring]
Put $z=b$; the altered team is $e_0=(0,1,z)$, of worth three.
The other $2b-1$ teams in $P$, called hard teams, have worth five.
Removing $e_0$ from the label--team incidence graph leaves $z$ isolated
and a tree on the other $4b-1$ labels: the remaining ring is a path,
and each $f_j$ attaches two leaves to a hub. Call this the main tree.
Root it at any label with any sector color. Traversing outward, color
each newly visited label by its sector position in the traversed team.
Every hard team then has at least two endpoints with their designated
colors, so every pair in it has at least one correctly colored endpoint.
Let $c^A$ be the coloring rooted at $0$ with color $A$, and $c^B$
the coloring rooted at $1$ with color $B$. The ring path gives
\[
                c^A(0)=c^A(1)=A,\qquad c^B(0)=c^B(1)=B.
\]

To attain ideal three at coordinate $(k,v)$ with $v\ne z$, root the
main tree at $v$ with color $k$, extend its coloring by $c(z)=C$,
and pay
\[
                        x_{s,j}=\one_{\{s=k\}}+2\one_{\{c(j)=s\}}.
\]
For a target at $z$, pay three to its target sector there and zero to
the other two sectors. On every main-tree label use, respectively,
\begin{align*}
 \text{target }A:\quad x_{s,j}&=\one_{\{s=A\}}
                         +\one_{\{c^A(j)=s\}}+\one_{\{c^B(j)=s\}},\\
 \text{target }B:\quad x_{s,j}&=\one_{\{s=B\}}+2\one_{\{c^B(j)=s\}},\\
 \text{target }C:\quad x_{s,j}&=\one_{\{s=C\}}+2\one_{\{c^A(j)=s\}}.
\end{align*}
Every label has total three, the target receives three, and every player
in the target sector receives at least one.

We check all team inequalities. A hard team receives at least four
coloring units and one base unit, covering five. Endpoints of a hard
pair receive at least two coloring units, covering its value at most
two. If a team containing that pair also has a diagonal pair, its three
labels lie on the main tree, so it also receives a base unit and covers
worth at most three.

The only positive soft pairs are $AC(0,z)=1$ and $BC(1,z)=2$.
For a target on the main tree, $C_z$ has two coloring units and every
full team has a base unit, so any team containing a soft pair receives
at least three. For target $C$ at $z$, $C_z=3$ gives the same coverage.
For target $B$ at $z$, labels $0,1,z$ all have vector $(0,3,0)$.
The soft $BC$ pair receives three; the soft $AC$ pair receives zero
but every $B$ coordinate is at least one, and if there is also a diagonal
pair its $B$ label is $0$ or $z$, where it receives three. For target
$A$ at $z$, labels $0,1$ have vector $(2,1,0)$ and $z$ has vector
$(3,0,0)$. The soft $AC$ pair receives two. The soft $BC$ pair
receives one and every $A$ coordinate is at least one; with an additional
diagonal pair, the $A$ label is $1$ or $z$ and the team receives at
least three. The soft team itself receives three in each case.

Two positive off-diagonal pairs still force a $P$ team. Every other
nondiagonal team has at most one such pair and at most one diagonal
pair, so the preceding checks exhaust these cases. Without off-diagonal
pairs its worth is at most one, covered by the target-sector lower
bound. Diagonal teams receive exactly three. Thus every displayed
witness is a total-$12b$ core point. The diagonal matching has this
value, and its tightness bounds every ideal by three. Separate witnesses
attain three at every coordinate.

For the exact subsidy, start with one per player and add two to the
$A$ player of each hard team. These players are distinct. Hard teams
receive five, the soft team receives three, and all other teams have
worth at most three. This covers all coalitions and the floors, with
total $12b+2(2b-1)$. Conversely, the disjoint teams in $P$ have total
worth $5(2b-1)+3=10b-2$, and their $6b$ unmatched players require one
each. The resulting lower bound on protected total is $16b-2$, matching
the construction.
\end{proof}

\begin{corollary}[No modulus uniform over market size]\label{anc:modulus}
Scale $R_b,R'_b$ by $1/(12b)$, giving grand worth one in both games,
and put $\delta_b=1/(6b)$. Their pair-matrix entrywise $\ell_1$
distance and their characteristic-function supremum distance both equal
$\delta_b$, whereas their protected subsidies differ by
$1/3-\delta_b$. Consequently no function $\omega(\delta)\to0$ as
$\delta\to0$ uniformly bounds this subsidy difference over all finite
market sizes in either distance.
\end{corollary}
\begin{proof}
Only one pair changes, by $\delta_b$. An integral matching uses that
pair at most once, so every coalition worth changes by an amount in
$[0,\delta_b]$. The three-player coalition $e_0$ changes by exactly
$\delta_b$. Positive scaling scales the core, ideals and subsidy,
so Theorem~\ref{anc:ring} gives the asserted subsidy difference.
\end{proof}
This is a failure of a size-independent modulus, not discontinuity at a
fixed market size. It also shows why removing a small pair value cannot
be charged a universal proportional cost when the original ideals change
with the market. Instability of selected assignment-core allocations
has established precedents~\cite[Example~3]{anc:KumabeYoshida2024}; the
scalar protected-subsidy statement and its priority require their own
comparison.

\subsection{An optimal obstruction that must contain cycles}
\begin{proposition}[A finite certificate]\label{anc:cycles}
There is a three-sector market with $m=11$ labels, grand worth $T=1$,
nonempty entire core and subsidy $76/303$, whose unique optimal grand
matching is diagonal. Every optimal team-form subsidy dual contains six
specified disjoint teams, each with multiplier one. Their connected
label--team incidence graph has cycle rank two. In particular, no
forest-supported optimal team-form dual exists, including fractional
duals or a change to another optimal grand matching.
\end{proposition}
The statement is computer-assisted through a supplied rational
certificate, not an asymptotic conclusion from experiments. We specify
the market and the finite proof obligations explicitly. All entries in
the following tables are divided by $5252$; every omitted pair entry
is zero. The diagonal $AB,AC,BC$ entries are
\[
\begin{array}{c|rrr@{\qquad}c|rrr}
 j&AB&AC&BC&j&AB&AC&BC\\\hline
 0&108&234&144&6&153&153&144\\
 1&252&144&90 &7&189&144&153\\
 2&144&198&198&8&144&144&198\\
 3&90&108&108 &9&198&144&198\\
 4&154&144&148&10&198&144&198\\
 5&108&234&144&&&&
\end{array}
\]
and the six off-diagonal teams $P$ have pair entries
\[
\begin{array}{c|rrr}
 (A,B,C)\text{ labels}&AB&AC&BC\\\hline
 (1,0,2)&252&288&252\\
 (3,4,1)&158&144&234\\
 (0,5,3)&252&198&108\\
 (6,3,7)&153&297&153\\
 (8,9,5)&288&234&306\\
 (7,10,8)&333&198&342
\end{array}
\]

\begin{proof}[Certificate proof]
The record \texttt{CYCLIC\_FRONTIER\_6\_2\_0} in the supplied file
\texttt{CYCLIC\_FRONTIER\_CERTIFICATES.json} contains all $33$ exact
ideals, a separate original-core point attaining each, their matching
upper certificates, and a protected allocation $y$ with a matching dual
of total $1+76/303$. Each core witness has nonnegative coordinates,
total one and all full-team inequalities. For the upper certificate of coordinate $i$, the stored multipliers
$\lambda_t\ge0$ and the efficiency multiplier $\nu$ satisfy
\[
 \nu-\sum_{t\ni j}\lambda_t\ge\one_{\{j=i\}}
 \quad(j\in N),\qquad
 \nu-\sum_t\lambda_tv(t)=H_i.
\]
Since $x\ge0$, $x(N)=1$ and $x(t)\ge v(t)$, these equations imply
$x_i\le\nu-\sum_t\lambda_tv(t)=H_i$. The multiplier on efficiency
is unrestricted; no unsupported joint ideal attainment is assumed. The protected vector satisfies all
floors and full-team inequalities, while the dual multipliers satisfy
the player-load bounds and the displayed objective. These exact finite
checks prove $T=1$, all declared ideals and the subsidy value; they do
not substitute a numerical solver tolerance for an inequality.

Average the $33$ original-core witnesses. Direct rational checks show
that the resulting point $x$ is strictly positive and its only tight
full teams are the eleven diagonals and $(1,0,2),(7,10,8)$.
A grand matching of worth one must cover every player and use only
$x$-tight teams: its value deficit is the sum of team slacks and
unmatched positive payoffs. At each label there is a player appearing
in no other tight team. Sectors
\[
                       A,B,A,A,A,A,A,B,A,A,A
\]
provide such a player at labels $0,\ldots,10$, respectively.
Their diagonal teams are all forced, proving uniqueness. This reasoning
also rules out another fractional perfect optimum.

For the protected problem use full-team multipliers $\lambda_t\ge0$
and floor multipliers $\mu_i\ge0$, with
$\sum_{t\ni i}\lambda_t+\mu_i\le1$. The primal--dual difference
at the certified $y$ is
\begin{align*}
 &\sum_t\lambda_t\bigl(y(t)-v(t)\bigr)
 +\sum_i\mu_i\bigl(y_i-H_i/3\bigr)\\
 &\hspace{25mm}+\sum_i y_i
                  \left(1-\sum_{t\ni i}\lambda_t-\mu_i\right).
\end{align*}
All terms are nonnegative. The certificate has $y_i>0$ for every
player, and exactly twelve full-team constraints are tight. The diagonal
teams through the six forcing players listed below are not among them;
this is included in the exact strict/tight checks. At any
optimal dual, non-tight teams have multiplier zero, strict-floor
players have $\mu_i=0$, and every player load is one. In order, the
following strict-floor players each occur in exactly one still
available tight team:
\[
\begin{array}{c|c}
 \text{player}&\text{forced team}\\ \hline
 A_0&(0,5,3)\\
 A_1&(1,0,2)\\
 B_9&(8,9,5)\\
 B_{10}&(7,10,8)\\
 C_1&(3,4,1)\\
 C_7&(6,3,7)
\end{array}
\]
Its multiplier is therefore one. All teams sharing any of its players
then have multiplier zero, justifying the exclusions before the next
step. The exact checker verifies each strict inequality and each unique
choice, rather than trusting the listed dual. Thus every optimal dual
contains all six teams with multiplier one. Their incidence graph is
connected with $18$ edges and $17$ vertices, so its cycle rank is
$18-17+1=2$. Adding further positive teams cannot remove these cycles.
\end{proof}
The certificate and checker are retained in
\nolinkurl{structural_classification_2026-09-06/sharp_subsidy/}.
The solver-free command
\texttt{python3 -B -O verify\_cyclic\_necessity.py}
reconstructs the market, verifies the original ideals and primal--dual
value, checks all strict/tight sets and forcing steps, and rejects a
false forcing step. Its direct dependency is
\texttt{verify\_cyclic\_frontier.py}, with its retained checker imports
and the named JSON certificate; run from the complete supplied directory.
This result rules out a forest-only optimal-dual restriction on the
same market. It does not exclude a transformation that changes the
market while correctly tracking all changes to its original ideals.

\paragraph{Reproducibility and claim boundaries.}
The same retained directory contains separate exact checkers for
Proposition~\ref{anc:compression}, Proposition~\ref{anc:attainment}
and Theorem~\ref{anc:ring}: respectively
\texttt{verify\_sector\_min\_core.py},
\texttt{verify\_coattainability.py}, and
\texttt{verify\_broken\_ring.py} together with \texttt{ring\_zero.py}.
They have been run with Python assertions disabled. The all-size
statements rely on the written proofs; the cyclic-optimum proposition
is expressly a finite certificate proof. Historical priority and
independent human specialist assessment remain open.
